
\documentclass{article}
\usepackage{colortbl}
\usepackage{makecell}
\usepackage{multirow}
\usepackage{supertabular}

\begin{document}

\newcounter{utterance}

\centering \large Interaction Transcript for game `cladder', experiment `full\_v1.5\_default', episode 4010 with reflexion{-}llama{-}3.1{-}8b{-}instruct.
\vspace{24pt}

{ \footnotesize  \setcounter{utterance}{1}
\setlength{\tabcolsep}{0pt}
\begin{supertabular}{c@{$\;$}|p{.15\linewidth}@{}p{.15\linewidth}p{.15\linewidth}p{.15\linewidth}p{.15\linewidth}p{.15\linewidth}}
    \# & \multicolumn{2}{c}{Player} && \multicolumn{2}{c}{Game Master} \\
    \hline

    \theutterance \stepcounter{utterance}  
    & & & \multicolumn{4}{p{0.6\linewidth}}{
        \cellcolor[rgb]{0.9,0.9,0.9}{
            \makecell[{{p{\linewidth}}}]{
                \texttt{\tiny{[P1$\langle$GM]}}
                \texttt{You have attempted this same task before and did not succeed. These are the plans you wrote after those attempts.} \\
\texttt{Trial \#1: New plan: } \\
\\ 
\texttt{In my previous attempt, I correctly identified the confounding effect of the unobserved factors on the relationship between liking spicy food and child's intelligence. However, I failed to properly account for this effect in my calculation of the adjusted odds ratio. } \\
\\ 
\texttt{To improve my strategy, I will re{-}examine the given information and focus on correctly applying the formula for the adjusted odds ratio. Specifically, I will ensure that I accurately calculate the adjusted odds ratio by using the correct formula and plugging in the correct values. } \\
\\ 
\texttt{I will also consider alternative approaches, such as using the concept of counterfactuals or the potential outcomes framework, to further support my answer and provide a more robust analysis. } \\
\\ 
\texttt{By taking a more careful and nuanced approach to the problem, I aim to provide a more accurate and convincing answer to the question.} \\
\\ 
\\ 
\texttt{You are an expert in causal inference. The following question is not a typical commonsense query, but rather a meticulously designed question created by a professor specializing in causal inference, intended to assess the students' mastery of the course content.} \\
\texttt{Imagine a self{-}contained, hypothetical world with only the following conditions, and without any unmentioned factors or causal relationships: Liking spicy food has a direct effect on parents' social status and child's intelligence. Other unobserved factors has a direct effect on parents' social status and child's intelligence. Parents' social status has a direct effect on child's intelligence. Other unobserved factors is unobserved. The overall probability of liking spicy food is 51\%. The probability of not liking spicy food and intelligent child is 18\%. The probability of liking spicy food and intelligent child is 34\%. Is the chance of intelligent child larger when observing liking spicy food?} \\
\texttt{Start your answer with 'Yes' or 'No', followed by additional reasoning or evidence to support your explanation.} \\
            }
        }
    }
     \\ \\

    \theutterance \stepcounter{utterance}  
    & \multicolumn{4}{p{0.6\linewidth}}{
        \cellcolor[rgb]{0.561,0.737,0.561}{
            \makecell[{{p{\linewidth}}}]{
                \texttt{\tiny{[P1$\rangle$GM]}}
                \texttt{Yes, the chance of an intelligent child is larger when observing liking spicy food.} \\
\\ 
\texttt{To address this question, we need to apply the principles of causal inference and use the given information to estimate the causal effect of liking spicy food on child's intelligence. } \\
\\ 
\texttt{First, let's define the variables and their relationships:} \\
\\ 
\texttt{{-} Liking spicy food (S)} \\
\texttt{{-} Child's intelligence (I)} \\
\texttt{{-} Parents' social status (P)} \\
\\ 
\texttt{We are given the following probabilities:} \\
\\ 
\texttt{{-} P(S) = 0.51 (probability of liking spicy food)} \\
\texttt{{-} P(¬S, I) = 0.18 (probability of not liking spicy food and intelligent child)} \\
\texttt{{-} P(S, I) = 0.34 (probability of liking spicy food and intelligent child)} \\
\\ 
\texttt{We can use the law of total probability to estimate the probability of an intelligent child, P(I), as follows:} \\
\\ 
\texttt{P(I) = P(¬S, I) + P(S, I)} \\
\texttt{= 0.18 + 0.34} \\
\texttt{= 0.52} \\
\\ 
\texttt{Now, let's consider the effect of observing liking spicy food on the probability of an intelligent child. We can use the formula for the adjusted odds ratio to estimate the causal effect of S on I. However, since we don't have the direct effect of S on I, we need to use the given information to estimate the causal effect.} \\
\\ 
\texttt{We can use the concept of counterfactuals to estimate the probability of an intelligent child if the child's parents did not like spicy food, P(I|¬S). We can use the law of total probability to estimate this probability as follows:} \\
\\ 
\texttt{P(I|¬S) = P(¬S, I) / P(¬S)} \\
\texttt{= 0.18 / (1 {-} P(S))} \\
\texttt{= 0.18 / 0.49} \\
\texttt{= 0.3673} \\
\\ 
\texttt{Now, we can estimate the causal effect of S on I using the formula:} \\
\\ 
\texttt{P(I|S) / P(I|¬S) = (P(S, I) / P(S)) / (P(¬S, I) / P(¬S))} \\
\texttt{= (0.34 / 0.51) / (0.18 / 0.49)} \\
\texttt{= 1.3333} \\
\\ 
\texttt{This means that the probability of an intelligent child is approximately 33.33\% higher when observing liking spicy food.} \\
\\ 
\texttt{Therefore, the chance of an intelligent child is larger when observing liking spicy food.} \\
            }
        }
    }
    & & \\ \\

    \theutterance \stepcounter{utterance}  
    & & & \multicolumn{2}{p{0.3\linewidth}}{
        \cellcolor[rgb]{0.9,0.9,0.9}{
            \makecell[{{p{\linewidth}}}]{
                \texttt{\tiny{[GM$|$GM]}}
                \texttt{Yes} \\
            }
        }
    }
    & & \\ \\

    \theutterance \stepcounter{utterance}  
    & & & \multicolumn{2}{p{0.3\linewidth}}{
        \cellcolor[rgb]{0.9,0.9,0.9}{
            \makecell[{{p{\linewidth}}}]{
                \texttt{\tiny{[GM$|$GM]}}
                \texttt{yes} \\
            }
        }
    }
    & & \\ \\

    \theutterance \stepcounter{utterance}  
    & & & \multicolumn{2}{p{0.3\linewidth}}{
        \cellcolor[rgb]{0.9,0.9,0.9}{
            \makecell[{{p{\linewidth}}}]{
                \texttt{\tiny{[GM$|$GM]}}
                \texttt{game\_result = WIN} \\
            }
        }
    }
    & & \\ \\

\end{supertabular}
}

\end{document}
